\chapter{Quantum observables under elimination}
\label{mht:ch:quantum-elimination}

The eliminated curve modes have a quadratic action on the Neumann
slice. Integrating them therefore determines both a determinant and
an inverse quadratic form. The determinant can cancel against the
ghost determinant while the inverse form still contracts observables.
This is the first quantum information that a classical residual
projection loses.

We begin with that Gaussian calculation, then retain its contractions
in the full observable map. Connected insertions give its higher
components. Background variation gives a connection on them.
The finite torus model and the first-jet boundary completion provide
quantum master equations on explicit finite spaces. Their proofs also
identify which additional constructions the smooth interacting theory
requires.

We then remove the curve and surface cutoffs on smooth
Gaussian observables, retaining their source-antifield corrections.
The continuum determinant uses an ordinary Sobolev space.

\section{The determinant and the covariance}

Use the cyclic Neumann algebra $E$, trace $\tau_E$, and residual
projections $p_0,p_1$ of Chapter~\ref{mht:ch:polar-neumann}.
Tensor with $\mathfrak g=\mathfrak{gl}_N(\mathbb C)$ and write
\[
 V=\ker p_0,\qquad W=\ker p_1,\qquad
 d_J=\bar\partial+\eta_0\operatorname{ad}_J.
\]
Here $V$ consists of matrix-valued functions and $W$ of matrix-valued
curve one-forms.

Theorem~\ref{mht:thm:background-contraction} supplies
the inverse $h_J:W\to V$. Put $h=h_0$ and
\[
 K_J(v,w)=\tau_E\operatorname{tr}(v d_Jw),\qquad
 M_J=h d_J|_V.
\]
The forms $K_J$ are skew. On the physical slice
$a=\eta_0J$, Proposition~\ref{mht:prop:exact-slice-action} identifies
the eliminated quadratic action with $K_J$ and the gauge derivative
with $M_J$.

Choose a finite-dimensional subspace $V_K\subset V$ on which $K_0$
is nondegenerate. Let $\pi_K$ be its $K_0$-orthogonal projection
and let $K_{J,K}$ be the matrix of $K_J|_{V_K}$ in one fixed basis.
Initially all expressions are formal in the entries of $J$ at zero.
The loop parameter $\hbar$ is independent of this completion.

\begin{proposition}\label{mht:prop:matched-gaussian}
The compressed gauge derivative is
\begin{equation}\label{mht:eq:matched-gauge-matrix}
 M_{J,K}=\pi_KM_J|_{V_K}=K_{0,K}^{-1}K_{J,K}.
\end{equation}

Normalize the bosonic action $y^{\mathsf T}K_{J,K}x$ relative to its
Gaussian at $J=0$. Normalize the ghost Gaussian with the same
matrix $M_{J,K}$. With the reference
orientations fixed, their vacuum factors are
\[
 Z_{\mathrm{bos}}=\det(M_{J,K})^{-1},\qquad
 Z_{\mathrm{gh}}=\det(M_{J,K}),\qquad
 Z_{\mathrm{bos}}Z_{\mathrm{gh}}=1.
\]
In the Wick convention with weight $e^{-y^{\mathsf T}K_{J,K}x/\hbar}$,
the normalized bosonic expectation satisfies
\begin{align}
 \langle x_i y_j\rangle
   &=\hbar(K_{J,K}^{-1})_{ij},\notag\\
 \langle x_i y_j x_k y_l\rangle
       -\langle x_i y_j\rangle\langle x_k y_l\rangle
   &=\hbar^2(K_{J,K}^{-1})_{il}(K_{J,K}^{-1})_{kj}.
 \label{mht:eq:matched-covariances}
\end{align}
\end{proposition}

\begin{proof}
Since $\bar\partial M_J=d_J$ on $V$, one has
$K_0(v,M_Jw)=K_J(v,w)$. Pairing with every $v\in V_K$ and using
the definition of $\pi_K$ proves
\eqref{mht:eq:matched-gauge-matrix}. Its constant term is the
identity. The block Hessian of the bosonic action has determinant
equal to a fixed sign times $\det(K_{J,K})^2$. The square-root branch
continued from the reference Gaussian gives the inverse determinant
ratio. Berezin integration gives the ghost determinant ratio.

For commuting source columns $F,G$, completing the square gives
\[
 \left\langle e^{F^{\mathsf T}x+G^{\mathsf T}y}\right\rangle
       =\exp\bigl(\hbar F^{\mathsf T}K_{J,K}^{-1}G\bigr).
\]
Two derivatives give the covariance. Four derivatives give two
pairings. Subtraction of the product of expectations leaves the
crossed pairing in \eqref{mht:eq:matched-covariances}.
\end{proof}

These formulas define a formal complex Gaussian, without a choice of
positive real contour. Numerical evaluation requires $K_{J,K}$ to
be invertible. The artificial singularity of
Example~\ref{mht:ex:compressed-pole} shows why invertibility of the
full smooth operator does not settle that requirement for an
arbitrary finite compression.

\begin{example}\label{mht:ex:native-gaussian-pair}
On the round sphere take
\[
 f=\frac{z^2}{(1+|z|^2)^2},\qquad
 g=\frac{\bar z}{1+|z|^2}.
\]
Both have zero $p_0$ projection. With $r=|z|^2$, polar integration
gives
\[
 \tau_E(f\bar\partial g)
       =-\kappa\int_0^\infty\frac{dr}{(1+r)^4}
       =-\frac\kappa3.
\]
Thus on their span
\[
 K_{0,K}=\kappa\begin{pmatrix}0&-1/3\\1/3&0\end{pmatrix},
 \qquad \langle x_1y_2\rangle=\frac{3\hbar}{\kappa}.
\]
The combined vacuum is one and this covariance is nonzero.

For a background-dependent example put
\[
 v=\frac r{1+r}-\frac12,\qquad
 u=\frac{zr}{(1+r)^2}.
\]
The two polar matrices on $(v,u)$ are
\[
 \kappa^{-1}\bigl(\tau_E(v_i\bar\partial v_j)\bigr)
    =\begin{pmatrix}0&1/6\\-1/6&0\end{pmatrix},\qquad
\kappa^{-1}\bigl(\tau_E(\eta_0v_iv_j)\bigr)
    =\begin{pmatrix}1/12&0\\0&0\end{pmatrix}.
\]

The nonzero entry of the second matrix is the integral
$\int_0^\infty(r/(1+r)-1/2)^2(1+r)^{-2}\,dr=1/12$.
The other entries vanish by angular weight.

Tensor with $\mathfrak{gl}_2$ and take $J=j\operatorname{diag}(1,-1)$.
In the order $E_{11},E_{12},E_{21},E_{22}$, let $G$ be the trace
matrix and let $L=\operatorname{ad}_J$. Direct multiplication gives
\begin{equation}\label{mht:eq:matrix-gaussian-inverse}
 \begin{gathered}
 M_{J,K}=\begin{pmatrix}1&0\\L/2&1\end{pmatrix},\qquad
 K_{0,K}=\kappa\begin{pmatrix}0&G/6\\-G/6&0\end{pmatrix},\\
 K_{J,K}^{-1}=\kappa^{-1}
       \begin{pmatrix}0&-6G^{-1}\\6G^{-1}&3LG^{-1}\end{pmatrix}.
 \end{gathered}
\end{equation}
Although $\det M_{J,K}=1$, the last block has entries
$6j/\kappa$ and $-6j/\kappa$. The covariance remembers the background.
\end{example}

The inverse quadratic form also has a smooth domain independent of
finite compression. Let $U\subset\mathbb R^2$ be an oriented surface
patch with coordinates $(s,t)$ and let $J:U\to\mathfrak g$ be smooth.
For smooth compactly supported maps $\phi,\gamma:U\to W$, define
\[
 F_\phi(x)=\int_U\tau_E\operatorname{tr}(x\phi)\,ds\,dt,
 \qquad
 G_\gamma(y)=\int_U\tau_E\operatorname{tr}(y\gamma)\,ds\,dt
\]
and
\[
 C_J(\phi,\gamma)=\int_U\tau_E\operatorname{tr}
                   \bigl((h_{J(s,t)}\gamma)\phi\bigr)\,ds\,dt.
\]
These covectors are represented by the cyclic pairing.

\begin{proposition}\label{mht:prop:smooth-gaussian-observables}
The generating rule
\[
 \mathcal P_J e^{F_\phi+G_\gamma}
       =e^{\hbar C_J(\phi,\gamma)}
\]
defines a normalized Gaussian functional on polynomials in the
displayed linear observables. Each Wick coefficient is continuous
on the smooth test spaces with a common compact surface support.
\end{proposition}

\begin{proof}
The equation $d_Jh_J\gamma=\gamma$ identifies $h_J\gamma$ with the
inverse kinetic form on these covectors. The background Green map
is continuous, with seminorm bounds uniform on compact matrix sets.
The image of $J$ on a common compact surface support is such a set.
Continuity of multiplication and the polar trace therefore proves
continuity of $C_J$. Every polynomial coefficient is a finite sum
of products of this bilinear form.
\end{proof}

The surface delta distribution has been paired with tests before
evaluation. The curve inverse still uses the global Green map.
Products of coincident local insertions require a further extension
of this domain.

\section{Finite BV integration}

Let $\mathcal F=\mathcal F'\oplus\mathcal F''$ be a finite-dimensional
graded symplectic complex with pairing of degree $-1$, split into
symplectic subcomplexes. Assume $\mathcal F''$ is acyclic. A cyclic
strong deformation retract onto $\mathcal F'$ fixes a gauge-fixing
Lagrangian in $\mathcal F''$ on which the free quadratic action is
nondegenerate. Fix compatible constant reference semidensities.
Integration on this Lagrangian, denoted $\mathcal I$, is the
coefficientwise Gaussian integral, including Berezin integration
in odd coordinates.

Write $\Delta,\Delta'$ for the degree-one BV Laplacians on the two
spaces. The odd bracket is normalized by
\begin{equation}\label{mht:eq:quantum-product-identity}
 \Delta(fg)=(\Delta f)g+(-1)^{|f|}f\Delta g
                              +(-1)^{|f|}\{f,g\}.
\end{equation}
The constant semidensity identifies functions with semidensities.
An even \emph{master weight} is an element $U$ of the integration
domain satisfying $\Delta U=0$. For $U=e^{I/\hbar}$ this is
\[
 \frac12\{I,I\}+\hbar\Delta I=0.
\]
The weight $e^{-S/\hbar}$ corresponds to $I=-S$ in this convention.

We use polynomial observables completed in $\hbar$ and in a positive
interaction parameter, with polynomial coefficients at each
interaction order. Alternatively, interactions of field degree at
least three use the completed field algebra
$\widehat{\operatorname{Sym}}(\mathcal F^\vee)[[\hbar]]$.
In the latter case the connected expansion is finite at each fixed
field degree and loop order. These are perturbative completions.
An analytic integration cycle requires its own convergence conditions.

\begin{proposition}\label{mht:prop:normalized-bv-pushforward}
Finite BV integration satisfies $\Delta'\mathcal I=\mathcal I\Delta$.
Let $I=I_0+I_{\mathrm{int}}$ be a degree-zero action with free
quadratic part $I_0$ and no negative powers of $\hbar$.
Require $I_{\mathrm{int}}$ to have positive interaction order in the
polynomial completion, or field degree at least three in the field
completion. For the master weight $U=e^{I/\hbar}$, assume
$U'=\mathcal I(U)$ has invertible leading coefficient after the
chosen Gaussian normalization, and put
\begin{equation}\label{mht:eq:normalized-bv-map}
 D=\hbar U^{-1}\Delta U,\qquad
 D'=\hbar(U')^{-1}\Delta'U',\qquad
 p(f)=(U')^{-1}\mathcal I(Uf).
\end{equation}
Then $D^2=(D')^2=0$, $D1=D'1=0$, $p1=1$, and $D'p=pD$.
Successive normalized integrations compose when each uses the
effective weight produced by the preceding integration.
\end{proposition}

\begin{proof}
The Laplacian splits into retained and eliminated terms. Retained
derivatives commute with the integral. Each eliminated term is an
integrated total derivative. It vanishes by Gaussian integration by
parts in even coordinates and by the Berezin derivative identity in
odd coordinates.

This proves the Stokes relation and implies $\Delta'U'=0$.
The exponential weights give
$D=\{I,-\}+\hbar\Delta$ and
$D'=\{I',-\}'+\hbar\Delta'$, where $U'=e^{I'/\hbar}$ is
interpreted in the same perturbative Gaussian domain. These
operators preserve the stated observable completions. Conjugating
the square-zero Laplacians proves the differential identities, and
\[
 D'p(f)=\hbar(U')^{-1}\Delta'\mathcal I(Uf)
       =\hbar(U')^{-1}\mathcal I\Delta(Uf)=pD(f).
\]
For two integrations put $U_1=\mathcal I_1U$ and
$U_2=\mathcal I_2U_1$, so the two intermediate factors cancel in
\[
 U_2^{-1}\mathcal I_2\bigl(U_1U_1^{-1}\mathcal I_1(Uf)\bigr)
       =U_2^{-1}\mathcal I_2\mathcal I_1(Uf).
\]
All identities are coefficientwise in the declared perturbative
domain.
\end{proof}

In \eqref{mht:eq:normalized-bv-map}, $U^{-1}\Delta U$ denotes
conjugation of operators. It is not multiplication by $U^{-1}\Delta(U)$.
For this finite symplectic splitting, the perturbative pushforward
is a quasi-isomorphism. Homological perturbation gives its inverse
and homotopy.\footnote{Cattaneo--Mnev,
\emph{BV pushforward as a quasi-isomorphism}, Theorem~5.1 and
Corollary~5.2, \href{https://arxiv.org/abs/2605.30558v2}{arXiv:2605.30558v2}.}
The eliminated acyclic complex and the coefficient completion are
part of this theorem's hypotheses.

For a free propagator $P$, include the factor $1/2$ in the contraction
operator $\partial_P$. Normalize its sign by the free Stokes identity
\[
 (Q'+\hbar\Delta')\,\pi e^{\hbar\partial_P}
       =\pi e^{\hbar\partial_P}(Q+\hbar\Delta),
\]
where $\pi$ restricts to retained fields and $Q,Q'$ are the free
differentials. It follows by commuting $Q$ through the exponential
and using the contraction identity.

If an interaction $I$ obeys
$(Q+\hbar\Delta)e^{I/\hbar}=0$, its effective interaction is
\begin{equation}\label{mht:eq:finite-effective-interaction}
 I'=\hbar\log\bigl(\pi e^{\hbar\partial_P}e^{I/\hbar}\bigr).
\end{equation}
The same intertwining proves the effective master equation.
The logarithm selects connected graphs. A graph with $v$ interaction
vertices and $e$ contractions has power $\hbar^{e-v+1}$, its loop
number, and coefficient $1/|\operatorname{Aut}\Gamma|$ times its
contracted vertex tensor. Thus \eqref{mht:eq:finite-effective-interaction}
constructs interaction vertices as well as the unary differential.

\begin{example}\label{mht:ex:scalar-effective-action}
Let $a\ne0$ and set
\[
 S(x,y)=\frac{a+\lambda y}{2}x^2+gxy^2+\frac b2y^2.
\]
Treat $g,\lambda$ formally. Completing the square and dividing the
$x$ integral by its reference Gaussian gives
\begin{equation}\label{mht:eq:scalar-effective-action}
 S_{\mathrm{eff}}(y)=\frac b2y^2
       -\frac{g^2y^4}{2(a+\lambda y)}
       +\frac\hbar2\log\frac{a+\lambda y}{a}.
\end{equation}

The conditional mean and variance are
\[
 p_y(x)=-\frac{gy^2}{a+\lambda y},\qquad
 p_y(x^2)-p_y(x)^2=\frac\hbar{a+\lambda y}.
\]

The rational correction in \eqref{mht:eq:scalar-effective-action}
is its tree series. The logarithm is its one-loop series. Evaluation
at the conditional mean recovers the former and loses the latter
and the variance.
\end{example}

\section{Connected insertions and their compositions}

Let $p:A\to B$ be the normalized map of
Proposition~\ref{mht:prop:normalized-bv-pushforward}. The observable
algebras $A,B$ are graded commutative and $\hbar$-adically complete.
Their differentials satisfy
\begin{equation}\label{mht:eq:observable-product-identity}
 D(fg)=(Df)g+(-1)^{|f|}fDg+(-1)^{|f|}\hbar\{f,g\},
\end{equation}
and the corresponding identity with primes on $B$.

For homogeneous inputs $f_1,\ldots,f_n$, let
$\operatorname{Part}(n)$ denote the partitions of $\{1,\ldots,n\}$.
Order each block by its input order and order blocks by their least
element. The sign $\epsilon(\pi;f)$ is the Koszul sign of this
reordering. Define the cumulant
\begin{equation}\label{mht:eq:quantum-cumulants}
 \kappa_n^p(f_1,\ldots,f_n)
  =\sum_{\pi\in\operatorname{Part}(n)}
     (-1)^{|\pi|-1}(|\pi|-1)!\,\epsilon(\pi;f)
       \prod_{C\in\pi}p\left(\prod_{i\in C}f_i\right).
\end{equation}

\begin{theorem}\label{mht:thm:quantum-cumulants}
In the finite perturbative Gaussian domain, with observables having
no negative powers of $\hbar$, one has
$\kappa_n^p\in\hbar^{n-1}B$. Hence the maps
\[
 \Phi_n^p=\hbar^{1-n}\kappa_n^p
\]
are defined over the original quantum power-series ring.
Their first component is $p$, and every higher component vanishes
when an input is the unit.

These maps give a coalgebra morphism for the degree-one coderivations
whose unary parts are $D,D'$ and whose binary parts are
$b_2(f,g)=(-1)^{|f|}\{f,g\}$. In particular,
\begin{align}
 D'\Phi_2^p(f,g)-\Phi_2^p(Df,g)
       -(-1)^{|f|}\Phi_2^p(f,Dg)
   =(-1)^{|f|}\bigl(p\{f,g\}-\{pf,pg\}'\bigr).
 \label{mht:eq:binary-quantum-ward}
\end{align}
For successive normalized integrations $p,q$, their full Taylor
series compose: $\Phi^{qp}=\Phi^q\circ\Phi^p$.
\end{theorem}

\begin{proof}
Introduce formal graded parameters $t_i$ with $|t_i|=-|f_i|$, so
$F=\sum_i t_i f_i$ has degree zero. The partition formula is
equivalent to
\begin{equation}\label{mht:eq:connected-observable-generator}
 \Phi^p(F)=\hbar\log p(e^{F/\hbar})
       =\sum_{n\ge1}\frac1{n!}\Phi_n^p(F,\ldots,F).
\end{equation}

The logarithm retains connected components containing all marked
observables. A connected graph with $n$ marks, $v$ interaction
vertices and $e$ contractions has $e\ge n+v-1$. Its quantum power
in the unscaled cumulant is $e-v$, proving divisibility by
$\hbar^{n-1}$. Normalization removes components containing only
interaction vertices.

At fixed interaction order the inputs are
polynomials, so the contributing contractions are finite. The
field-completed alternative is finite at each specified field degree
and loop order.

For completeness, let $S^c(A)=\bigoplus_{n\ge1}S^n(A)$ be the
conilpotent symmetric coalgebra. The coalgebra automorphism $\tau_A$
whose Taylor coefficients are the products in $A$ sums over
partitions. Its inverse is finite on each tensor length.

Since $D$ has order at most two, its conjugated coderivation has only
unary part $D$ and binary part $\hbar b_2$.
The conjugate of $S^c(p)$ has Taylor coefficients
\eqref{mht:eq:quantum-cumulants}. The relation $D'p=pD$ gives its
coalgebra identity. Dividing the arity-$n$ identity by
$\hbar^{n-1}$ gives the displayed normalized coderivations and
their morphism. Direct subtraction of
\eqref{mht:eq:observable-product-identity} before and after $p$
also proves \eqref{mht:eq:binary-quantum-ward}.

For a scalar mark $t1$, the generating function factors by
$e^{t/\hbar}$. Its logarithm has no mixed coefficient with this
mark, proving the unit assertion. Finally
\[
 e^{\Phi^p(F)/\hbar}=p(e^{F/\hbar})
\]
and substitution into the next logarithm proves composition,
including every partition coefficient and Koszul sign.
\end{proof}

The quantum division is performed before specializing to a quotient
by $\hbar^r$. Knowing an unscaled cumulant only modulo $\hbar^r$
does not determine its divided component to that same precision.

\begin{example}\label{mht:ex:gaussian-cumulants}
For one even Gaussian coordinate $x$ with invertible stiffness $A$,
\[
 p_A(x^{2m})=(2m-1)!!(\hbar/A)^m,\qquad
 p_A(x^{2m+1})=0.
\]

Its BV complex adjoins $\eta$ of degree $-1$, with
$D(\eta f)=Axf-\hbar f'$ and $D(f)=0$. Gaussian integration by
parts proves its Stokes identity. The first nonzero connected maps are
\begin{equation}\label{mht:eq:composite-gaussian-cumulants}
 \Phi_2(x,x)=A^{-1},\qquad
 \Phi_3(x^2,x,x)=2A^{-2},\qquad
 \Phi_n(x^2,\ldots,x^2)=2^{n-1}(n-1)!\hbar A^{-n}.
\end{equation}
The second formula is
$\hbar^{-2}(p_A(x^4)-p_A(x^2)^2)$.
For the last, expand
$\log p_A(e^{t x^2})=-\tfrac12\log(1-2\hbar t/A)$.
Thus the higher comparison remains nontrivial even when $p_A(x)=0$.
\end{example}

\section{Changing the covariance and lifting sources}

Let $C^{ab}(J)$ be a symmetric matrix of formal functions of even
parameters $J$. On polynomials in even coordinates $x^a$, define
\[
 p_J=\operatorname{ev}_{x=0}
       \exp\left(\frac\hbar2 C^{ab}(J)\partial_a\partial_b\right).
\]
Repeated indices are summed. On each polynomial the exponential
contains finitely many contractions. Write $d_{\mathrm{par}}$ for
exterior differentiation in the parameters.

\begin{proposition}\label{mht:prop:covariance-connection}
The connection
\begin{equation}\label{mht:eq:covariance-connection}
 \nabla^{\mathrm W}=d_{\mathrm{par}}+
       \frac\hbar2(d_{\mathrm{par}}C^{ab})\partial_a\partial_b
\end{equation}
is flat and obeys $d_{\mathrm{par}}(p_JF)=p_J\nabla^{\mathrm W}F$.
Its transport from $J_0$ to $J_1$ is
\begin{equation}\label{mht:eq:covariance-transport}
 F_{J_1}=\exp\left[-\frac\hbar2
          \bigl(C(J_1)-C(J_0)\bigr)^{ab}\partial_a\partial_b\right]
          F_{J_0},\qquad p_{J_1}F_{J_1}=p_{J_0}F_{J_0}.
\end{equation}
\end{proposition}

\begin{proof}
All constant-coefficient contraction operators commute. Differentiate
their exponential to obtain \eqref{mht:eq:covariance-connection}.
The exterior derivative of its one-form is zero and its commutator
square is zero, proving flatness. Multiplication of the two commuting
exponentials proves \eqref{mht:eq:covariance-transport}.
\end{proof}

For $C(J)=(1+J)^{-1}$, the quadratic expectation has derivative
$\partial_Jp_J(x^2)=-\hbar/(1+J)^2$.
The connection accounts for this term. A varying field frame also
has the cotangent source shift and the semidensity factor of
\eqref{mht:eq:moving-antifield-shift} and
\eqref{mht:eq:moving-density-action}. These data enter the same
observable comparison, with different roles.

A finite gauge quartet displays the source shift directly.
Let $J$ be even, with conjugate $J^+$ of degree $-1$. Take even
coordinates $x,b$ and odd coordinates $c,\bar c$ of degrees $1,-1$.
For vector-valued coordinates all following expressions sum over
components. Set
\[
 Q_0x=c,\qquad Q_0\bar c=b,\qquad Q_0b=Q_0c=Q_0J=0.
\]

Adjoin cotangent antifields $q^+$ of degree $-1-|q|$. Fix
$\{q,q^+\}=1$, use left odd derivatives, and set
\[
 \Delta=\sum_q(-1)^{|q|}\partial_q\partial_{q^+},\qquad
 S_0=\sum_i(x_i^+c_i-\bar c_i^+b_i).
\]
This convention gives $\{S_0,q\}=Q_0q$. For an invertible formal
matrix $M(J)$ with $M(0)=1$, let
\[
 \Psi=\bar cM(J)x,\qquad
 Q_0\Psi=bM(J)x-\bar cM(J)c.
\]

\begin{proposition}\label{mht:prop:quantum-source-lift}
The canonical transformation $T=\exp(-\{\Psi,-\})$ has unit
Berezinian, commutes with $\Delta$, and sends $S_0$ to
$S_\Psi=S_0+Q_0\Psi$. Both actions satisfy the quantum master
equation. For $D_\Psi=\{S_\Psi,-\}+\hbar\Delta$, the transformed
source
\begin{equation}\label{mht:eq:quantum-source-lift}
 T(J^+)=J^+-\bar c M'(J)x
\end{equation}
is closed. For several parameters the corresponding term is
$\bar c(\partial_{J_a}M)x$, and $T$ transports every polynomial
in their conjugates.
\end{proposition}

\begin{proof}
The cotangent lift of $Q_0$ squares to zero and has zero divergence,
so $\{S_0,S_0\}=\Delta S_0=0$. The fermion has no antifields.
Thus $\Delta\Psi=0$ and its Hamiltonian flow commutes with $\Delta$.
The flow fixes fields and translates antifields, giving a triangular
Jacobian with identity diagonal. The transformation of $S_0$
terminates at its first commutator because $Q_0\Psi$ has no
antifields. This proves the action and master equations.

The derivative of $Q_0\Psi$ with respect to $J$ is
$bM'x-\bar cM'c$. It equals $Q_0(\bar cM'x)$, whereas
$\{S_\Psi,J^+\}=\partial_J(Q_0\Psi)$. Thus the Hamiltonian
differential kills \eqref{mht:eq:quantum-source-lift}. Its Laplacian
is zero as well. More generally, $T$ intertwines the two quantum
differentials and is an algebra map. The original polynomials in
the $J_a^+$ are closed, proving the last assertion.
\end{proof}

\section{A homotopy for a changing master weight}

The source correction is a special case of a normalized variation.
Keep a finite BV integration map $\mathcal I$ and the Laplacians
fixed. Let $t$ be a parameter and let $R_t$ have degree $-1$.
Assume a family of master weights satisfies
\begin{equation}\label{mht:eq:master-weight-family}
 \partial_tU=\Delta(R_tU),\qquad \Delta U=0.
\end{equation}
Use the normalized maps of \eqref{mht:eq:normalized-bv-map} and put
$r_t=p_t(R_t)$. Stokes gives
$\partial_tU'=\Delta'(r_tU')$.

For multiplication $m_R$, set $L_R=[\Delta,m_R]$, the graded
commutator. Define degree-zero operators
\[
 v_A=U^{-1}L_RU-m_{U^{-1}\partial_tU},\qquad
 v_B=(U')^{-1}[\Delta',m_r]U'
                         -m_{(U')^{-1}\partial_tU'}.
\]
An order-two operator commuted with multiplication has order at
most one. Subtracting its value at $1$ therefore makes each $v$
a derivation. Moreover
$\partial_tD=[v_A,D]$ and $\partial_tD'=[v_B,D']$, by
\eqref{mht:eq:master-weight-family} and $\Delta^2=(\Delta')^2=0$.

\begin{proposition}\label{mht:prop:normalized-gauge-homotopy}
The degree-minus-one map
\[
 h_t(f)=\hbar^{-1}\bigl(p_t(R_tf)-p_t(R_t)p_t(f)\bigr)
\]
satisfies
\begin{equation}\label{mht:eq:normalized-gauge-homotopy}
 \partial_tp_t+p_tv_A-v_Bp_t=D'h_t+h_tD,\qquad h_t1=0.
\end{equation}
If $R_t$ has no negative quantum powers in the perturbative Gaussian
domain, the division by $\hbar$ is integral.
\end{proposition}

\begin{proof}
On unnormalized weights put $K=\mathcal I m_R-m_r\mathcal I$.
The Stokes relation gives
\[
 \mathcal I L_R-L_r\mathcal I=\Delta'K+K\Delta.
\]

Differentiate $p_t=(U')^{-1}\mathcal I U$. In
$\partial_tp_t+p_tv_A-v_Bp_t$ the two logarithmic derivatives
cancel, leaving the conjugate of the left side. The right side
becomes $D'h_t+h_tD$ because
$h_t=\hbar^{-1}(U')^{-1}KU$. Normalization gives $K(U)=0$.
The numerator of $h_t$ is a two-input cumulant, so the integrality
follows from Theorem~\ref{mht:thm:quantum-cumulants}.
\end{proof}

The homotopy retains higher insertions as well. Write $F_p$ for the
undivided cumulant coalgebra map, and $\mathcal Q_A,\mathcal Q_B$
for the coderivations with unary parts $D,D'$ and binary parts
$\hbar b_2$.

Let $S^c(p_t;h_t)$ apply $h_t$ to one factor and $p_t$ to every
other factor, summed with the graded tensor signs. Define
\[
 \mathcal H_t=\tau_B^{-1}S^c(p_t;h_t)\tau_A.
\]
Derivations commute with the multiplication coalgebra maps.
Conjugating
\eqref{mht:eq:normalized-gauge-homotopy} gives
\[
 \partial_tF_p+F_p\widetilde v_A-\widetilde v_BF_p
       =\mathcal Q_B\mathcal H_t+\mathcal H_t\mathcal Q_A,
\]
where $\widetilde v$ is the unary coderivation induced by $v$.
The arity-$n$ coefficient of $\mathcal H_t$ is
\begin{equation}\label{mht:eq:marked-quantum-homotopy}
 (\mathcal H_t)_n(f_1,\ldots,f_n)
        =\hbar^{-1}\kappa_{n+1}^{p_t}(R_t,f_1,\ldots,f_n).
\end{equation}

Indeed, differentiating a partition term replaces one block
expectation by its connected expectation with $R_t$. The first
summands enumerate partitions in which the block of $R_t$ is not
a singleton. For a partition into $k$ other blocks, the $k$
subtracted summands contribute $(-1)^k k!$, the coefficient of
the additional singleton block. This proves the formula, with
$R_t$ placed first to fix the odd signs.

Connected graphs now have $n+1$ marks, so
$(\mathcal H_t)_n\in\hbar^{n-1}B$. Its divided coefficients give
the homotopy of the normalized higher maps. For two successive
eliminations, the undivided coalgebra formula is
\[
 \mathcal H^{qp}=F_q\mathcal H^p+\mathcal H^qF_p,
\]
where $F_p,F_q$ are the cumulant coalgebra maps. This follows by
differentiating $F_{qp}=F_qF_p$, and is associative under three
eliminations.

A higher parameter simplex requires a total master weight satisfying
its parameter-dependent master equation. Separate paths alone do not
provide that datum.

\section{Finite surface actions and the torus}

Let $\ell_n$ be the cyclic $L_\infty$ operations on
$E_K\otimes\mathfrak{gl}_N$ obtained in
Theorem~\ref{mht:thm:finite-cyclic-transfer}. Their degrees are $2-n$.
For an oriented surface $U$, tensor with $\Omega^\bullet(U)$,
placing surface forms first, and include $d_U$ in the unary
differential. Write $\langle x,y\rangle=\tau_E\operatorname{tr}(xy)$
for the internal pairing, extended with the tensor signs.

For a total-degree-one superfield $\mathsf A$, define
\begin{equation}\label{mht:eq:finite-surface-action}
 S_{K,U}=\frac12\int_U\langle\mathsf A,(d_U+d_E)\mathsf A\rangle
     +\sum_{n\ge2}\frac1{(n+1)!}
          \int_U\langle\mathsf A,\ell_n(\mathsf A^{\otimes n})\rangle.
\end{equation}
The series is completed in field amplitude. Its odd symplectic form
and cohomological vector field are
\[
 \omega_{K,U}=\frac12\int_U\langle\delta\mathsf A,\delta\mathsf A\rangle,
 \qquad
 Q\mathsf A=(d_U+d_E)\mathsf A
       +\sum_{n\ge2}\frac1{n!}\ell_n(\mathsf A^{\otimes n}).
\]

\begin{proposition}\label{mht:prop:finite-surface-boundary}
For closed $U$, or compact support in its interior,
\eqref{mht:eq:finite-surface-action} satisfies the classical master
equation. If $U$ has boundary, its cohomological vector field and
field-space symplectic form satisfy
\[
 \iota_Q\omega_{K,U}=\delta S_{K,U}+\pi^*\alpha_{K,\partial U},
 \qquad
 \alpha_{K,\partial U}
      =\frac12\int_{\partial U}\langle\mathsf A,\delta\mathsf A\rangle,
\]
where $\pi$ is boundary restriction. Internal transfer preserves
surface support and commutes with restriction. The boundary terms
cancel on an interface with opposite induced orientations.
\end{proposition}

\begin{proof}
Cyclicity makes the variation of each $(n+1)$-field term equal to
its Hamiltonian contraction with $\ell_n/n!$. The transferred
identities give $Q^2=0$. Integration by parts in $d_U$ leaves the
displayed boundary primitive. The other operations use only internal
coefficients and commute with $d_U$ and restriction, with their
Koszul signs. Stokes' theorem supplies the opposite interface signs.
\end{proof}

To obtain a finite quantum model, replace the surface complex by
its specified Frobenius algebra
\[
 B=H^*(T^2;\mathbb C)=\Lambda(a,b),\qquad
 |a|=|b|=1,\qquad \tau_B(ab)=1.
\]
Its differential and its other basis traces are zero. The shifted
field space
\[
 \mathcal F_K=(B\otimes E_K\otimes\mathfrak{gl}_N)[1]
\]
has dimension $(8+16K)N^2$. The pairing
$\tau_B\tau_E\operatorname{tr}$ induces an odd symplectic form.
Fix a constant Darboux semidensity, its Laplacian $\Delta_K$, and
the action $S_K$ obtained from \eqref{mht:eq:finite-surface-action}
by replacing integration over $U$ with $\tau_B$.

\begin{theorem}\label{mht:thm:finite-torus-qme}
The action obeys
\[
 \frac12\{S_K,S_K\}+\hbar\Delta_KS_K=0.
\]

For $K\le L$, normalized finite BV integration of the acyclic
complement of $\mathcal F_K$ in $\mathcal F_L$ gives $S_K$ with
no loop correction to its action. These integrations compose.
\end{theorem}

\begin{proof}
For a homogeneous Frobenius basis $(e_i)$, use right trace duals
$\tau_B(e_je^i)=\delta_j^i$. Closing two inputs of a multiplication
tensor inserts the Euler element
\[
 \mathfrak e_B=\sum_i(-1)^{|e_i|}e_ie^i.
\]
The basis $(1,a,b,ab)$ has dual $(ab,b,-a,1)$, so
\begin{equation}\label{mht:eq:torus-euler-element}
 \mathfrak e_B=ab-ab-ab+ab=0.
\end{equation}

Each vertex of $S_K$ factors into a multiplication tensor of $B$
and an internal cyclic tensor. Applying $\Delta_K$ closes two
inputs and therefore gives zero. The classical master equation
follows from the transferred identities, proving the quantum one.

The eliminated propagator is the identity on $B$ tensored with
the internal complement homotopy. In a connected interaction graph,
contract a spanning tree on the $B$ factors. The Frobenius dual-basis
identity joins the multiplication tensors into one. Every remaining
edge inserts $\mathfrak e_B$. Thus every graph with a loop vanishes.

The surviving trees are exactly the nested transfer of
\eqref{mht:eq:finite-nested-homotopy}. Normalizing the free Gaussian
to one and using coefficientwise integration proves composition.
\end{proof}

Open contractions survive, since $\tau_B(ab)=1$. They enter the
cumulant maps of Theorem~\ref{mht:thm:quantum-cumulants}.
For a closed genus-$g$ surface, the Euler element of its cohomology
algebra is $(2-2g)$ times the normalized top class. Thus the loop
cancellation uses the chosen torus coefficient algebra.

The full surface heat propagator imposes a second restriction.
On a flat patch, let $L_a,L_b$ denote wedge multiplication by the
parallel one-forms and let $\iota_a,\iota_b$ be the dual exterior
contractions. Pure wedge multiplication on $\Lambda(a,b)$ has
supertrace zero: positive-degree maps are strictly upper triangular,
and the identity has supertrace $1-2+1=0$.

Exterior contractions change this trace:
\begin{equation}\label{mht:eq:surface-loop-detector}
 \operatorname{Str}(L_aL_b\iota_b\iota_a)=1.
\end{equation}
Indeed, in the order $1,a,b,ab$ this operator is the projection
onto $ab$. The adjoint surface differential in the full heat
propagator contains such exterior contractions. The Euler argument
therefore does not cancel every coefficient tensor of its loops.
The spacetime kernels and the full graph sum are still needed to
compute an anomaly.

\section{The compensating first-jet field}

The first-jet boundary in \eqref{mht:eq:neumann-pullback}
permits an explicit enlargement of the gauge group. Work in the
formal normalized-frame neighborhood of the polar connection $a$.
Its frame $H(a)$ obeys $\bar\partial H=-aH$, $H(o)=I_N$, and its
residual is $\mathcal J(a)=-j(H(a))$.
For a formal smooth gauge $k$ whose amplitude-zero coefficient is
an invertible constant matrix, put
$k_0=k(o)$ and $k_1=j(k)$, without imposing $k_1=0$.
Frame uniqueness gives
$H(a^k)=k^{-1}H(a)k_0$, and hence
\begin{equation}\label{mht:eq:affine-residual-gauge}
 \mathcal J(a^k)=k_0^{-1}\mathcal J(a)k_0+k_0^{-1}k_1.
\end{equation}
Here the derivative of $k^{-1}$ is
$-k_0^{-1}k_1k_0^{-1}$, fixing the affine sign.

Introduce a degree-zero edge matrix $\phi$ with the same affine
law. Then
\[
 X=\mathcal J(a)-\phi,\qquad
 X^k=k_0^{-1}Xk_0,\qquad
 \mathcal O_r^{\mathrm{rel}}=\frac1r\operatorname{tr}(X^r)
 \quad(r\ge1).
\]
For formal $\phi$ of positive amplitude, the smooth global gauge
$k=e^{-f_0\phi}$ has $k_0=I_N$, $k_1=-\phi$ and sends $\phi$ to
zero. The subgroup preserving $\phi=0$ is exactly the Neumann
group. The edge field compensates the translation part of
$\operatorname{GL}_N(\mathbb C[\xi]/\xi^2)$. It is different from
the coefficient $a_1$ of the boundary superconnection in
\eqref{mht:eq:polar-bfv-bf}.

The finite residual complex has even matrices $J,\phi$ and odd
ghost matrices $c_0,c_1$, of degrees $1,1$. Its algebra is polynomial
in even entries and exterior in odd entries. Define the degree-one
derivation
\begin{align}
 QJ&=c_1+[J,c_0],& Q\phi&=c_1+[\phi,c_0],\notag\\
 Qc_0&=-c_0^2,& Qc_1&=-c_0c_1-c_1c_0.
 \label{mht:eq:first-jet-brst}
\end{align}
Matrix multiplication includes the exterior signs of ghost entries.
The comparison with smooth fields sets $J=\mathcal J(a)$.

\begin{theorem}\label{mht:thm:edge-brst-contraction}
The derivation $Q$ squares to zero. In the polynomial coordinates
\[
 X=J-\phi,\qquad \gamma=c_1+[\phi,c_0]
\]
it is
\[
 QX=[X,c_0],\quad Qc_0=-c_0^2,\quad Q\phi=\gamma,\quad Q\gamma=0.
\]

Let $p$ kill $\phi,\gamma$, let $i$ include the retained
Chevalley--Eilenberg complex
\[
 C_N=\mathbb C[X_{ij}]\otimes\Lambda((c_0)_{ij}),
\]
and let $N_{\mathrm{pair}}$ count the added variables. On positive
pair degree set
\[
 h=N_{\mathrm{pair}}^{-1}\sum_{ij}
                         \phi_{ij}\partial_{\gamma_{ij}},
\]
with $h=0$ in degree zero. Then $Qh+hQ=1-ip$.
\end{theorem}

\begin{proof}
The graded Leibniz rule verifies $Q^2=0$ on the four matrix
generators. It gives the displayed equations after the coordinate
change, whose inverse is $J=X+\phi$,
$c_1=\gamma-[\phi,c_0]$.
With left odd derivatives, the anticommutator of
$\sum\gamma_{ij}\partial_{\phi_{ij}}$ and
$\sum\phi_{ij}\partial_{\gamma_{ij}}$ is
$N_{\mathrm{pair}}$. The retained differential involves only
$X,c_0$ and graded-commutes with the added homotopy. Division on
each positive eigenspace proves the identity.
\end{proof}

The retained complex contains its full ghost cohomology. Replacing
it by invariant functions of $X$ would discard part of this
retraction. The first invariant has the particularly simple form
\begin{equation}\label{mht:eq:relative-linear-source}
 \mathcal O_1^{\mathrm{rel}}
       =\kappa^{-1}\tau_E(\operatorname{tr}a)-\operatorname{tr}\phi.
\end{equation}
To see this, apply polar Stokes to
$\bar\partial\log\det H=-\operatorname{tr}a$ and use
$j(\log\det H)=-\operatorname{tr}\mathcal J(a)$.
Both terms on the right of \eqref{mht:eq:relative-linear-source}
have variation $\operatorname{tr}c_1$ under the enlarged gauge
group, so their difference is closed.

\section{The finite edge master equation and local primitives}

Adjoin antifields $q^+$ of degree $-1-|q|$ to every coordinate
of \eqref{mht:eq:first-jet-brst}. Use the cotangent bracket and
Laplacian of Proposition~\ref{mht:prop:quantum-source-lift}. The
cotangent Hamiltonian is
\[
 S_Q=\sum_q(-1)^{|q|}q^+Qq,\qquad \{S_Q,q\}=Qq.
\]

\begin{theorem}\label{mht:thm:finite-edge-qme}
The coordinate change of Theorem~\ref{mht:thm:edge-brst-contraction}
and its cotangent lift have Berezinian one. They split $S_Q$ into
the retained cotangent Hamiltonian and
\[
 S_{\mathrm{pair}}=\sum_{ij}\phi_{ij}^+\gamma_{ij}.
\]
Both $\{S_Q,S_Q\}$ and $\Delta S_Q$ vanish. For every
conjugation-invariant polynomial or formal field-only potential
$V(X)$,
\begin{equation}\label{mht:eq:finite-edge-qme}
 \frac12\{S_Q+V,S_Q+V\}+\hbar\Delta(S_Q+V)=0.
\end{equation}
Normalized BV integration of the separated added pair leaves
the retained potential unchanged.
\end{theorem}

\begin{proof}
The even change $J=X+\phi$ is a determinant-one shear. The odd
change $c_1=\gamma-[\phi,c_0]$ is triangular with identity diagonal.
Their base Berezinian, and hence that of the cotangent lift, is one.
The split differential determines $S_{\mathrm{pair}}$ and also gives
$Q\gamma^+=\phi^+$, $Q\phi^+=0$.

The cotangent master bracket is the lift of $2Q^2$, so vanishes.
For the Laplacian, the adjoint action on $\mathfrak{gl}_N$ has
trace zero. The first-jet Lie algebra has adjoint matrix
\[
 \begin{pmatrix}\operatorname{ad}_x&0\\
                 \operatorname{ad}_y&\operatorname{ad}_x\end{pmatrix}
\]
at $x+\xi y$, and therefore also has trace zero. These identities
give zero divergence for the even representation and ghost parts
of \eqref{mht:eq:first-jet-brst}. The affine terms have zero
divergence. Thus $\Delta S_Q=0$.

The potential has no antifields, so $\Delta V=\{V,V\}=0$.
Invariance gives $\{S_Q,V\}=0$, completing the master equation.

On the added Lagrangian $\phi=\gamma^+=0$, the remaining
coordinates $\phi^+,\gamma$ are odd and have the nondegenerate
quadratic pairing $\phi^+\gamma$. Orient each Berezin integral by
$\int\phi^+\gamma=1$ and normalize it by $\hbar$. Then
$\hbar\int e^{\phi^+\gamma/\hbar}=1$.
The product over entries is one and the retained weight factors out.
\end{proof}

The contraction also gives supported primitives in a precisely
defined local algebra. On a surface $U$, take polynomial densities
in finitely many derivatives of the retained fields, the added
fields and all their antifields. For a derivative multi-index $I$,
write $\phi_I,\gamma_I,\phi_I^+,\gamma_I^+$ for the added jets.
Prolongation of the differential gives
\[
 Q\phi_I=\gamma_I,\quad Q\gamma_I=0,\quad
 Q\gamma_I^+=\phi_I^+,\quad Q\phi_I^+=0.
\]
Let $N_{\mathrm{add}}$ count these four sorts of variables, and let
$p$ kill them. Define
\[
 h_0=\sum_{I,i,j}
   \left(\phi_{I,ij}\partial_{\gamma_{I,ij}}
          +\gamma_{I,ij}^+\partial_{\phi_{I,ij}^+}\right),\qquad
 h=N_{\mathrm{add}}^{-1}h_0
\]
on positive added degree, and set $h=0$ on degree zero.

\begin{theorem}\label{mht:thm:supported-edge-primitive}
On these densities and their quotient by total surface derivatives,
one has $Qh+hQ=1-ip$. The maps preserve surface support and maximal
derivative order. They act coefficientwise on formal series of
finite jet order in each coefficient. They also extend to a fixed
normal-current coefficient complex with defined products and total
Koszul signs.
Consequently a defined local residual $\mathcal A$ satisfying
$Q\mathcal A=0$ and $p\mathcal A=0$ has primitive
\begin{equation}\label{mht:eq:supported-edge-primitive}
 \mathcal A=Q(h\mathcal A).
\end{equation}
\end{theorem}

\begin{proof}
Each field pair and each antifield pair contributes its Euler
operator to $Qh_0+h_0Q$. The retained differential uses only retained
jets, so graded-commutes with $h_0$. Division by the positive count
proves the identity.

A total derivative shifts each pair to the
next jet pair and commutes with $h_0$ and the count. The identity
therefore descends modulo total derivatives. No inverse differential
operator or new derivative is introduced. Coefficient functions
remain at the same surface point. This proves the support and
derivative-order assertions.

Tensoring with the specified coefficient
complex preserves the identity. Apply it to the closed residual
with zero retained projection to obtain
\eqref{mht:eq:supported-edge-primitive}.
\end{proof}

For a compactly supported smooth density $\rho$,
\[
 \int_U\rho\,\operatorname{tr}c_1
       =Q\left(\int_U\rho\,\operatorname{tr}\phi\right).
\]
The same homotopy leaves a retained component $p\mathcal A$
untouched. Its coefficient and cohomology class must be computed
in the retained theory.

\section{A retained current anomaly}

The first-jet boundary has an algebraic deformation that makes this
last restriction concrete. Let
\[
 R=\mathbb C[\sigma,\nu],\qquad
 B_\sigma=R[\zeta]/(\zeta^2-\sigma),\qquad
 \tau(1)=0,\quad\tau(\zeta)=\kappa\in\mathbb C^\times.
\]
Here $\sigma,\nu$ have degree zero and $\zeta$ is a coefficient
coordinate, distinct from a surface or current coordinate. The
fiber at $\sigma=0$ is the first-jet Frobenius algebra.

Let
$\mathfrak g$ be a finite-dimensional complex Lie algebra with
nondegenerate invariant symmetric form $\kappa_0$, and set
$\mathfrak d_\sigma=\mathfrak g\otimes B_\sigma$.
Writing $(x,v)=x+\zeta v$, its bracket and form are
\begin{align}
 [(x,v),(y,w)]&=([x,y]+\sigma[v,w],\ [x,w]+[v,y]),\notag\\
 B_0((x,v),(y,w))&=\kappa\bigl(\kappa_0(x,w)+\kappa_0(v,y)\bigr).
 \label{mht:eq:collision-current-form}
\end{align}
The subalgebra $\mathfrak l=R\otimes_{\mathbb C}\mathfrak g$,
embedded through the coefficient $1$, is Lagrangian.
Denote the Killing form by
$K_{\mathfrak g}(x,y)=\operatorname{Tr}_{\mathfrak g}
(\operatorname{ad}_x\operatorname{ad}_y)$.

For an invariant level $k$, the affine current algebra has modes
\[
 [J_m(X),J_n(Y)]=J_{m+n}([X,Y])
               +m\delta_{m+n,0}k(X,Y)\mathbf1.
\]

Induce its vacuum module from the one-dimensional module on which
all nonnegative modes vanish and the central element acts as one.
Ordered negative modes give its PBW basis.
Choose a basis $t_a$ of $\mathfrak l$ with structure constants
$f_{ab}{}^c$, and odd ghosts $c^a,b_a$ of degrees $1,-1$ with
$b_a(w)c^b(w')\sim\delta_a^b/(w-w')$.
The current coordinate is $w$ and the BRST residue is
\begin{equation}\label{mht:eq:affine-brst-residue}
 Q_{\mathrm{cur}}=\operatorname{Res}_w
       \left({:}c^aJ_a{:}
          -\frac12 f_{ab}{}^c{:}c^ac^bb_c{:}\right).
\end{equation}

\begin{proposition}\label{mht:prop:collision-current-level}
The invariant level
\begin{equation}\label{mht:eq:collision-current-level}
 k_\sigma=\nu B_0-\frac12K_{\mathfrak d_\sigma}
    =\begin{pmatrix}-K_{\mathfrak g}&\nu\kappa\kappa_0\\
                    \nu\kappa\kappa_0&-\sigma K_{\mathfrak g}
      \end{pmatrix}
\end{equation}
defines a universal affine vacuum algebra free over $R$.
On its tensor product with the ghosts,
$Q_{\mathrm{cur}}^2=0$ over $R$ and in every fiber.
After $\sigma=a^2$ and inversion of $a$, the two separated levels are
\[
 k_\pm=\pm\frac{\nu\kappa}{2a}\kappa_0
                         -\frac12K_{\mathfrak g}.
\]
\end{proposition}

\begin{proof}
The adjoint matrix of $(x,v)$ is
\[
 \begin{pmatrix}\operatorname{ad}_x&\sigma\operatorname{ad}_v\\
                 \operatorname{ad}_v&\operatorname{ad}_x\end{pmatrix}.
\]
The trace of a product gives
$K_{\mathfrak d_\sigma}((x,v),(y,w))
=2K_{\mathfrak g}(x,y)+2\sigma K_{\mathfrak g}(v,w)$.
This proves \eqref{mht:eq:collision-current-level} and its restriction
$k_\sigma|_{\mathfrak l}=-K_{\mathfrak g}$. Invariance proves the
Jacobi identity of the mode algebra. PBW and induction give the
free vacuum module without inverting either parameter.

Compute the two residues of the product of
\eqref{mht:eq:affine-brst-residue} with itself. The current double
pole contributes
$\tfrac12k|_{\mathfrak l,ab}\operatorname{Res}{:}\partial_wc^ac^b{:}$.
The terms containing a current and two ghosts cancel against the
mixed contractions. The remaining single contractions cancel by
the Lie Jacobi identity.

The two ghost contractions give the
adjoint trace, with contribution
$\tfrac12K_{\mathfrak g,ab}\operatorname{Res}{:}\partial_wc^ac^b{:}$.
Thus, in these conventions,
\begin{equation}\label{mht:eq:retained-brst-square}
 Q_{\mathrm{cur}}^2=\frac12
       (k|_{\mathfrak l}+K_{\mathfrak g})_{ab}
                     \operatorname{Res}{:}\partial_wc^ac^b{:}.
\end{equation}
The chosen level makes this zero.
Finally the idempotents $(1\pm\zeta/a)/2$ have traces
$\pm\kappa/(2a)$ and separate the two adjoint representations.
They give the stated levels on the split fiber.
\end{proof}

At the bare level $\nu B_0$, the restriction to $\mathfrak l$
vanishes and leaves the ghost term. For $\mathfrak g=\mathfrak{sl}_2$
with $[h,e]=2e$, $[h,f]=-2f$, one has $K_{\mathfrak g}(h,h)=8$.
Set $c^h=\eta_+w+\eta_-w^{-1}$ with independent odd coefficients.
Then
\[
 \operatorname{Res}(\partial_wc^h)c^h=2\eta_+\eta_-,\qquad
 Q_{\mathrm{cur}}^2=8\eta_+\eta_-\ne0.
\]

This polynomial in retained ghosts survives the projection that
kills the added edge pairs. The zero-projection hypothesis of
Theorem~\ref{mht:thm:supported-edge-primitive} therefore fails.
The half-Killing level supplies the algebraic correction.

Its polynomial dependence on $\nu$ does not identify $\nu$ with
an inverse physical coupling at zero coupling. That normalization
still requires the native field-theory comparison.

\section{Smooth Gaussian Ward identities}

The smooth functional of
Proposition~\ref{mht:prop:smooth-gaussian-observables} is defined
without a cutoff. We now obtain it from finite cyclic spaces, with
control of background derivatives. This comparison concerns the
quadratic complement in a prescribed smooth residual field.

Use the staggered spaces $\widehat V_L$, projections
$\widehat\Pi_L$, and operators $X_L$ of
Chapter~\ref{mht:ch:angular-cutoffs}. Put
\[
 M_{J,L}=1+X_L\otimes\operatorname{ad}_J,\qquad
 h_{J,L}=M_{J,L}^{-1}\widehat\Pi_Lh.
\]
The inverse exists for every finite matrix $J$. If
$a\in\widehat V_L\otimes\mathfrak g$ and
$\gamma\in W$, then
\begin{equation}\label{mht:eq:staggered-pairing-inverse}
 K_{J,L}(a,h_{J,L}\gamma)=\tau_E\operatorname{tr}(a\gamma).
\end{equation}

Indeed, $K_{J,L}(a,v)=K_0(a,M_Jv)$ and
$\widehat\Pi_LM_Jh_{J,L}\gamma=\widehat\Pi_Lh\gamma$.
Adjointness of the projection for $K_0$ proves the identity.
Thus $h_{J,L}$ inverts the restricted kinetic form on the smooth
covectors represented by this pairing.

For the compactly supported tests $\phi,\gamma$ used above, define
\[
 C_{J,L}(\phi,\gamma)=\int_U\tau_E\operatorname{tr}
                  \bigl((h_{J(s,t),L}\gamma)\phi\bigr)\,ds\,dt.
\]
All tests have common compact surface support. Backgrounds and
variations range over bounded smooth sets, with the background
values in a fixed compact subset of $\mathfrak g$.

\begin{theorem}\label{mht:thm:smooth-ward-limit}
At every fixed polynomial degree, the normalized Wick functionals
\[
 \mathcal P_{J,L}e^{F_\phi+G_\gamma}
       =e^{\hbar C_{J,L}(\phi,\gamma)}
\]
converge to $\mathcal P_J$, rapidly in $L$ on the stated bounded
sets. The convergence holds after every fixed finite number of
background variations and surface derivatives.
For fixed smooth matrix variations $A,B$,
\begin{align}
 \partial_Ah_J&=-h_J\eta_0\operatorname{ad}_A h_J,\notag\\
 \partial_B\partial_Ah_J
   &=h_J\eta_0\operatorname{ad}_B h_J\eta_0\operatorname{ad}_A h_J
     +h_J\eta_0\operatorname{ad}_A h_J\eta_0\operatorname{ad}_B h_J.
 \label{mht:eq:ordered-smooth-responses}
\end{align}
The same identities hold for $h_{J,L}$ and converge to these
formulas. The matrices $A,B,J$ need not commute.
\end{theorem}

\begin{proof}
Apply Theorem~\ref{mht:thm:uniform-staggered-inverse} to $h\gamma$.
Its adapted Hilbert norms control ordinary Sobolev norms, and the
polar pairing is continuous on the smooth core. For $s\geq3$,
Sobolev multiplication and $\tau_E(\eta)=\kappa j(G\eta)$ give
\[
 |\tau_E\operatorname{tr}(v\phi)|
       \leq C_s\|v\|_{H^s}\|\phi\|_{H^s}.
\]

Integration over the fixed compact support preserves the estimate.
Surface derivatives act on the tests and on the finite matrix
parameters of the inverse. The theorem already controls every
fixed number of those derivatives. Each Wick coefficient is a
finite sum of products of covariances, so it has the same rapid
convergence, with a larger seminorm bound.

Differentiating $M_JM_J^{-1}=1$ gives
\[
 \partial_A M_J^{-1}
       =-M_J^{-1}(X\otimes\operatorname{ad}_A)M_J^{-1}.
\]
Use $X=h\eta_0$ and $h_J=M_J^{-1}h$ to obtain the first response.
A second differentiation gives both ordered terms in
\eqref{mht:eq:ordered-smooth-responses}. The finite projection is
independent of $J$, so the same calculation applies at finite $L$.
The parameter-derivative estimates identify their limits.
\end{proof}

For a finite list of smooth smearings, write $x_a,y_b$ for the
resulting even linear observables and $C_{ab}(J)$ for their
covariance without $\hbar$. On their polynomial algebra,
\[
 \mathcal P_J=\operatorname{ev}_0
       \exp\left(\hbar\sum_{a,b}C_{ab}(J)
                           \partial_{x_a}\partial_{y_b}\right).
\]
The corresponding Ward connection is
\begin{equation}\label{mht:eq:smooth-ward-connection}
 \nabla_A^{\mathrm W}=\partial_A+
       \hbar\sum_{a,b}(\partial_AC_{ab})
                           \partial_{x_a}\partial_{y_b},
 \qquad
 \partial_A(\mathcal P_JF)=\mathcal P_J\nabla_A^{\mathrm W}F.
\end{equation}

Mixed parameter derivatives and the constant-coefficient
contraction operators commute. The connection is therefore flat
for fixed parameter directions, as in
Proposition~\ref{mht:prop:covariance-connection}. It is compatible
with changing or enlarging the list of smearings, so defines a
linear connection on the smooth cylindrical polynomial core.
Its second-order term retains the contractions in the observable
map. The connection is not a derivation of the ordinary product.

\subsection{Surface cells and source antifields}

The quadratic complement has no surface derivatives. This permits
an independent finite surface regulator. Choose a bounded oriented
coordinate rectangle $R\Subset U$ containing the support of the
tests, and partition it into cells $Q_\nu$ of diameter at most $a$.
Write $w_\nu=\int_{Q_\nu}d\sigma>0$ and choose
$\sigma_\nu\in Q_\nu$.

At each cell take the quartet of
Proposition~\ref{mht:prop:quantum-source-lift}. Its even fields are
$x_\nu,b_\nu$, and its odd fields $c_\nu,\bar c_\nu$ have degrees
$1,-1$. Their components
lie in the finite curve space and its dual. The reference kinetic
pairing identifies $b_\nu$ with the second surface component.
The quadratic action and gauge fermion are
\begin{align}
 S_{J,L,a}&=\sum_\nu w_\nu
       \bigl(b_\nu M_{J(\sigma_\nu),L}x_\nu
             -\bar c_\nu M_{J(\sigma_\nu),L}c_\nu\bigr),\notag\\
 \Psi_{J,L,a}&=\sum_\nu w_\nu
                \bar c_\nu M_{J(\sigma_\nu),L}x_\nu.
 \label{mht:eq:surface-quartet-regulator}
\end{align}

Use ordinary Darboux brackets in each cell and the reference
Berezin orientation for which, with $M=M_{J(\sigma_\nu),L}$,
\[
 \mathcal P(x_{\nu,i}b_{\nu,j})
       =\hbar w_\nu^{-1}(M^{-1})_{ij},\qquad
 \mathcal P(\bar c_\nu Bc_\nu)
       =\hbar w_\nu^{-1}\operatorname{tr}(BM^{-1}).
\]
Different cells have no contractions. Smooth smearings use the
weights $w_\nu$, so their covariance is the Riemann sum for
$C_{J,L}$.

Before evaluating the residual matrices at $J(\sigma_\nu)$,
give each $J_\nu$ its ordinary cotangent antifield $J_\nu^+$.
For a prescribed smooth direction $A$, put
\[
 J_A^+=\sum_\nu\langle A(\sigma_\nu),J_\nu^+\rangle,
 \qquad A_{\nu,L}=X_L\otimes\operatorname{ad}_{A(\sigma_\nu)}.
\]
The gauge fermion has the closed source lift
\begin{equation}\label{mht:eq:surface-quartet-source}
 \mathcal L_{A,L,a}=J_A^+
       -\sum_\nu w_\nu\bar c_\nu A_{\nu,L}x_\nu.
\end{equation}
The same cell weights occur in the action, fermion and source
correction. Closure follows by the canonical transformation of
Proposition~\ref{mht:prop:quantum-source-lift} before any limit.

For a fixed finite parameter family $J+\sum_i z_iA_i$, the same
formula uses the retained antifield $z_i^+$ in place of $J_{A_i}^+$.
We keep this residual coefficient algebra fixed as the cells are
refined.

The quartet cylindrical core consists of polynomials in smoothly
smeared $x,b,c,\bar c$ and unintegrated residual parameters and
their antifields. Eliminated cotangent antifields and unsmeared
coincident surface composites are excluded. Let
$\mathcal P_J^{\mathrm q}$ be the Wick functional on this core
with continuum covariance $C_J$, the displayed ghost signs, and
the identity on residual variables. It extends $\mathcal P_J$.
All formal moments are interpreted coefficientwise in $\hbar$.

\begin{proposition}\label{mht:prop:joint-quartet-limit}
On every fixed cylindrical polynomial, the finite functionals
$\mathcal P_{J,L,a}$ converge to $\mathcal P_J^{\mathrm q}$ as
$L\to\infty$ and $a\to0$. For fixed polynomial degree, every
integer $r\geq0$ and bounded smooth sets, the error is at most
\[
 C_r\bigl((1+L)^{-r}+a\bigr).
\]
The estimate holds after any fixed finite number of background
variations. Moments containing finitely many source primitives
\eqref{mht:eq:surface-quartet-source} also have limits, with the
same estimate. The finite Ward insertions define convergent
operations on this core.
\end{proposition}

\begin{proof}
For a smooth function $f$ on $R$, the mean value estimate gives
\[
 \left|\sum_\nu w_\nu f(\sigma_\nu)-\int_Rf\,d\sigma\right|
       \leq a\,\operatorname{Area}(R)\sup_R|\nabla f|.
\]
Theorem~\ref{mht:thm:smooth-ward-limit} bounds the required
derivatives of the finite covariance integrands uniformly in $L$.
Its rapid curve estimate and this quadrature bound prove the
claim for covariances. Fixed Wick polynomials have finitely many
contractions, proving the first assertion and its differentiated
versions.

Each correction $\bar cA_{\nu,L}x$ has no internal contraction,
and two such corrections cannot contract with each other.
In a moment against a fixed cylindrical polynomial, both ends
must therefore attach to external fields. The resulting open
chains involve the bounded direction operator
$\widehat\Pi_L(X\otimes\operatorname{ad}_A)\widehat\Pi_L$
and the convergent inverse. Each source primitive contributes
one cell weight. The contractions leave the weights of an
ordinary Riemann sum. This proves the assertion for their
products.

For Ward insertions, differentiate the actual finite Gaussian.
Its normalized determinant is one in every cell. For a test $F$
independent of $J$, let $W_{A,L,a}=\partial_AS_{J,L,a}$. Then
\begin{align}
 \mathcal P_{J,L,a}(W_{A,L,a})&=0,\notag\\
 \partial_A\mathcal P_{J,L,a}(F)
       &=-\hbar^{-1}\mathcal P_{J,L,a}(W_{A,L,a}F),\notag\\
 \partial_B\partial_A\mathcal P_{J,L,a}(F)
       &=\hbar^{-2}\mathcal P_{J,L,a}(W_{B,L,a}W_{A,L,a}F).
 \label{mht:eq:regulated-inserted-ward}
\end{align}

The last identity uses the linear dependence of $M_{J,L}$ on $J$.
The differentiated left sides converge with the stated estimate,
so the finite insertions on the right have those limits. This
constructs the inserted operations without multiplying surface
delta distributions after the limit.
\end{proof}

This regulator is local on the surface for the displayed quadratic
action, and global on the curve. Extending it to nonlinear surface
vertices requires their differential and local products as further
data.

\subsection{A source response with unit vacuum}

The first four-mode space gives a nonzero inserted response.
Use $(u,v,\mathcal Cv,\mathcal Cu)$ from
\eqref{mht:eq:angular-first-background}, tensor with
$(E_{11},E_{12},E_{21},E_{22})$, and put
\[
 H=\operatorname{diag}(1,-1),\qquad
 A=X_1\otimes\operatorname{ad}_H,\qquad
 M=1+\lambda A,\qquad
 C=M^{-1}=1-\lambda A+\lambda^2A^2-\lambda^3A^3.
\]
Here $A^4=0$ and $\lambda$ is an even residual parameter.
For the quartet action $S=bMx-\bar cMc$, normalize the Wick
functional $\mathcal P$ with weight $e^{-S/\hbar}$ and the
Berezin orientation used above. Adjoin cotangent antifields and
write
\[
 S_{\mathrm{cot}}=\sum_i(x_i^+c_i-\bar c_i^+b_i),\qquad
 D^- =\{S_{\mathrm{cot}}+S,-\}-\hbar\Delta.
\]
The source lift is
$\mathcal L=\lambda^+-\bar cAx$ and satisfies $D^-\mathcal L=0$.

Let $B$ take $v\otimes E_{12}$ to $u\otimes E_{12}$ and vanish
on the other basis vectors. Set $F=\bar cBc$ and
$W=bAx-\bar cAc$.

\begin{proposition}\label{mht:prop:inserted-source-detector}
The normalized moments satisfy
\begin{equation}\label{mht:eq:inserted-source-detector}
 \mathcal P(F)=\hbar\lambda,\qquad
 \mathcal P(W)=0,\qquad
 \mathcal P(WF)=-\hbar^2,\qquad
 \partial_\lambda\mathcal P(F)
       =-\hbar^{-1}\mathcal P(WF)=\hbar.
\end{equation}
The bare residual antifield obeys $D^-\lambda^+=W$, detected by
$\mathcal P((D^-\lambda^+)F)=-\hbar^2$. The corrected lift is
closed before integration.
\end{proposition}

\begin{proof}
The four-mode chain gives
\[
 \operatorname{tr}(AC)=0,\qquad
 \operatorname{tr}(BC)=\lambda,\qquad
 \operatorname{tr}(ACBC)=-1.
\]
Berezin contraction, with the order of odd factors fixed, gives
\[
 \mathcal P((\bar cAc)(\bar cBc))
   =\hbar^2\bigl(\operatorname{tr}(AC)\operatorname{tr}(BC)
                         -\operatorname{tr}(ACBC)\bigr).
\]
The independent bosonic factor gives
$\mathcal P((bAx)F)=\hbar^2\operatorname{tr}(AC)
\operatorname{tr}(BC)$. Subtraction proves the moment formulas.
Finally $D^-\lambda^+=W=D^-(\bar cAx)$ by differentiation of
the gauge fermion, proving both closure statements.
\end{proof}

This is a test of the inserted cochain map. The ghosts remain
contractible. It persists on the smooth cylindrical core for a
constant background $J=\lambda H$: smear the two ghost factors
by smooth surface tests whose pairing is one. The weight change
from $u$ to $v$ permits exactly one $X$ step, with the same
coefficient in every larger staggered space. Thus the nonzero
response survives Proposition~\ref{mht:prop:joint-quartet-limit},
although the vacuum determinant is one.

\section{The regularized continuum determinant}

The adapted Hilbert space controls the actual staggered inverse.
For a determinant, use instead the ordinary constrained Sobolev
space
\[
 V_s^{\mathrm{ord}}
   =\{f\in H^{s+1}(C):j(f)=p_0(f)=0\},\qquad s\geq2.
\]
The operators agree on their common smooth vectors. On this
space $X=h\eta_0$ gains one ordinary derivative.
Fix a positive Hilbert norm on $\mathfrak g$ and put
$A_J=X\otimes\operatorname{ad}_J$, so $M_J=1+A_J$.

The Schatten ideal $\mathcal S_p$ consists of compact operators
whose singular values have summable $p$th powers. The ideal
$\mathcal S_1$ is the trace class. For an integer $p\geq1$ and
$A\in\mathcal S_p$, the regularized Fredholm determinant is
\begin{equation}\label{mht:eq:regularized-neumann-determinant}
 \det\nolimits_p(1+A)
  =\det\left((1+A)
       \exp\left(\sum_{r=1}^{p-1}\frac{(-1)^r}{r}A^r\right)\right).
\end{equation}
The operator inside the determinant is identity plus trace class.
The determinant on the right is the ordinary Fredholm determinant.

\begin{theorem}\label{mht:thm:neumann-continuum-determinant}
On $V_s^{\mathrm{ord}}$, the operator $X$ belongs to
$\mathcal S_p$ for every real $p>2$. For every integer $n\geq3$,
$A_J^n$ is trace class and $\operatorname{Tr}(A_J^n)=0$.
For every integer $p\geq3$ and every finite matrix $J$,
\[
 \det\nolimits_p M_J=1.
\]
The trace-log series and all its finite matrix-parameter
derivatives converge locally uniformly in $J$, in trace norm
before taking the trace.
\end{theorem}

\begin{proof}
The spherical Sobolev inclusion $H^{s+2}(C)\hookrightarrow
H^{s+1}(C)$ has singular values
$(1+\ell(\ell+1))^{-1/2}$ with multiplicity $2\ell+1$.
Their $p$th powers are summable for $p>2$, since the tail is
bounded by a multiple of $\sum_\ell\ell^{1-p}$.

Factor $X$ through this inclusion. The Green operator gains
one derivative, and its finite-rank corrections obey the same
bound because first-jet evaluation is continuous at these indices.
Fixed smooth corrections impose $j=p_0=0$. The ideal property
therefore gives $X\in\mathcal S_p$ for $p>2$.

The Schatten product inequality makes $X^3$ trace class, and
bounded composition gives the same for every $X^n$, $n\geq3$.
\footnote{B.~Simon, \emph{Notes on infinite determinants of Hilbert
space operators}. Advances in Mathematics \textbf{24} (1977),
244--273, \href{https://math.caltech.edu/SimonPapers/74.pdf}{\S\S1, 3, 6}.}

Rotation is unitary on the ordinary constrained space and
$U_\theta XU_\theta^{-1}=e^{i\theta}X$. Trace invariance gives
\[
 \operatorname{Tr}(X^n)=e^{in\theta}\operatorname{Tr}(X^n)=0
 \qquad(n\geq3),
\]
where the last equality follows by choosing $e^{in\theta}\ne1$.
The finite matrix factor can now be traced:
\[
 \operatorname{Tr}(A_J^n)=\operatorname{Tr}(X^n)
                    \operatorname{tr}_{\mathfrak g}
                         ((\operatorname{ad}_J)^n)=0.
\]

On this ordinary carrier $X$ is compact. A nonzero spectral
value would bring an entire nonzero rotation circle into its
spectrum, contradicting the compact-operator spectrum theorem.
Thus $\|X^n\|^{1/n}\to0$. For $n\geq p\geq3$,
\[
 \|X^n\|_1\leq\|X^p\|_1\|X^{n-p}\|.
\]
The root test gives absolute trace-norm convergence, locally
uniformly in $J$, of
\begin{equation}\label{mht:eq:neumann-trace-log}
 \sum_{n\geq p}\frac{(-1)^{n+1}}{n}A_J^n.
\end{equation}

A fixed number $k$ of matrix derivatives introduces at most
$n(n-1)\cdots(n-k+1)$ ordered placements of directions. This
polynomial factor leaves the root test unchanged.

The scalar series defining the operator in
\eqref{mht:eq:regularized-neumann-determinant} differs from one
first at order $p$. Factoring out $A^p$ proves its trace-class
property. Near $J=0$, its logarithmic determinant is the trace
of \eqref{mht:eq:neumann-trace-log}, and hence zero.
Analyticity of the regularized Fredholm determinant gives the
identity for every $J$.
\end{proof}

For a matrix direction $U$, set $A_U=X\otimes\operatorname{ad}_U$.
Differentiating the trace-log, using cyclicity only for trace-class
products, gives
\begin{equation}\label{mht:eq:neumann-determinant-connection}
 d\log\det\nolimits_p M_J(U)
   =(-1)^{p+1}\operatorname{Tr}
                   (M_J^{-1}A_J^{p-1}A_U)=0.
\end{equation}
Indeed, the derivative series is
$\sum_{n\geq p}(-1)^{n+1}\operatorname{Tr}(A_J^{n-1}A_U)$.
Summing its geometric factor proves the formula. Each summand
has positive rotation weight $n$ and zero trace. The same
trace-norm estimates justify every further finite source
derivative, retaining all ordered placements of its matrix
directions.

This determinant is a scalar entire function on the declared
Hilbert family. It assigns no ordinary trace to $X$ or $X^2$,
and does not define the full native quantum measure. The
nonzero response in \eqref{mht:eq:inserted-source-detector}
remains part of the observable comparison.

\subsection{The two subtracted orders}

A finite rotation-equivariant projector $\Pi$ gives
$\operatorname{Tr}((\Pi A_J\Pi)^n)=0$ for every $n\geq1$ on
its finite range. A trace-class regulator $R_\epsilon$ commuting
with rotations likewise gives
$\operatorname{Tr}(R_\epsilon A_J^n)=0$ for $n\geq1$ on
$V_s^{\mathrm{ord}}\otimes\mathfrak g$. These prescriptions assign zero to
the first two traces. Relaxing that normalization leaves the
following precise freedom for constant backgrounds.

\begin{proposition}\label{mht:prop:determinant-quadratic-normalization}
Let $N\geq2$ and let $F(J)$ be a holomorphic germ vanishing at zero.
Suppose its Taylor coefficients vanish in degrees at least three,
and it is invariant under conjugation and central translation
$J\mapsto J+cI_N$. Then, for a unique scalar $\beta$,
\begin{equation}\label{mht:eq:determinant-quadratic-normalization}
 F(J)=\beta\operatorname{tr}_{\mathfrak g}
                         ((\operatorname{ad}_J)^2)
     =\beta\bigl(2N\operatorname{tr}(J^2)
                           -2(\operatorname{tr}J)^2\bigr).
\end{equation}
If $F(e^{i\theta}J)=F(J)$ also holds, then $\beta=0$.
For $N=1$, central-translation invariance already gives $F=0$.
\end{proposition}

\begin{proof}
The Taylor hypothesis makes $F$ a polynomial of degree at most
two. On diagonal matrices it is symmetric. Matrices with distinct
eigenvalues are dense, so conjugation invariance gives
\[
 F=a\operatorname{tr}J+b\operatorname{tr}(J^2)
                               +c(\operatorname{tr}J)^2.
\]

Central translation forces $a=0$ and $c=-b/N$.
On $\operatorname{End}(\mathbb C^N)$ one has
$\operatorname{ad}_J=J\otimes1-1\otimes J^{\mathsf T}$.
Taking the trace of its square proves the displayed formula,
with $\beta=b/(2N)$. It is nonzero for some traceless $J$ when
$N\geq2$, proving uniqueness. Rotation invariance kills this
homogeneous term of degree two. In rank one every matrix is
a central translate of zero.
\end{proof}

The hypotheses compare constant-background logarithms with
$\log\det_3M_J=0$. They do not classify local cyclic counterterms
depending on surface derivatives, normal sources or boundary
fields.

\section{The interacting comparison}

The preceding constructions supply finite master weights, their
connected observables, supported contractions of an added edge
sector, and smooth Gaussian source limits. The regularized scalar
determinant and the inserted Ward responses have separate roles.
On the smooth theory the next equation involves the actual
regulated interacting fields. For a positive-scale kernel, an
effective interaction $I_t=\sum_{g\ge0}\hbar^g I_{g,t}$ must satisfy
\[
 QI_t+\frac12\{I_t,I_t\}_t+\hbar\Delta_tI_t=0.
\]

After lower orders have been solved, the next residual is closed
for their linearized differential. Its primitive must lie in the
prescribed local cyclic counterterm complex, with the same normal
support and boundary data.

\begin{problem}[The native quantum comparison]
\label{mht:qn:native-quantum-elimination}
Extend the smooth Gaussian regulator and source-antifield map to
one regulator for the interacting Neumann theory and its nonminimal
fields. Require compatibility with cyclic transfer and the polar
trace, estimates uniform on compact residual-matrix sets, and
Ward identities on the enlarged observable domain. Construct local
cyclic extensions at every simultaneous source collision.
Compute the retained interacting master residual and determine
whether it has an allowed supported primitive.

Construct the map from independently specified dynamical closed
fields and descendants to these observables, preserving the normal
marking, relative boundary operations and connected insertions.
Compare its determinant, residual measure, boundary polarization,
disk normalization and heat operator with an independently prepared
state, including sewing. In the affine collision family, derive the
level and its parameter normalization from the native regulator.
\end{problem}

The free kernels of Chapter~\ref{mht:ch:heat-contacts} determine the
positive-scale contractions. The finite torus theorem determines
the action after internal elimination on its specified coefficient
algebra. The remaining state comparison requires a preparation on
both sides of the elimination map.
